\documentclass{article}
\usepackage{ITB_MSCe}
\usepackage{tcolorbox}
\tcbuselibrary{breakable}
\usepackage{bold-extra} % bold tt
\usepackage{lipsum}
\usepackage{graphicx}
\usepackage{color}
\usepackage{amsthm}
\usepackage{amssymb}
\usepackage{amsmath}
\usepackage{amsfonts}
\usepackage{amstext}
\usepackage{latexsym}
\usepackage{upgreek}
\usepackage{hyperref}
\usepackage{url}
\usepackage{nicefrac}
\usepackage{pifont}
\usepackage{marvosym}
\usepackage{enumitem}
\usepackage{tikz}
\usepackage[numbers]{natbib}
\usepackage{eso-pic}
\usepackage{textcomp}
\usepackage{fancyvrb}
\usepackage{mwe}
\usepackage{lastpage}
\usepackage{listings, xcolor}
\usepackage{algorithmic}
\usepackage{algorithm}

\usepackage{array}
\usepackage{ragged2e}
%\bibpunct{}{}{;}{a}{,}{,}

\input{images/rgb}

\lstdefinestyle{python_code}{%
  language=Python, basicstyle=\ttfamily, columns=fullflexible,
  keywordstyle=\color{blue}, commentstyle=\itshape\color{green!50!black},
  stringstyle=\color{magenta}, showstringspaces=false,
  morekeywords={as,@staticmethod}, }

\lstset{ style=python_code }

\newsavebox{\sclbox}

\newcolumntype{L}[1]{>{\RaggedRight}p{#1}}



\newenvironment{declaration}
  {\noindent\textbf{Declaration:}} {\par}

\title{Explaining the policy--practice gap in AI education governance using
multi-layered NLP: A comparative study}

\author{
\normalsize
\begin{tabular}[t]{c@{\extracolsep{1em}}c} Mariem NSIRI & Geraldine GRAY\\
Technological University Dublin & Technological University Dublin\\
B00177462@myTUDublin.ie & supervisor.name@tudublin.ie
\end{tabular}
}

%\date{\today}

\begin{document}

\section{Methodology}

This study adopts a multi-layered text-as-data methodology to examine the
policy--practice gap in AI education governance. The purpose is not to evaluate
NLP methods for their own sake, but to use computational text analysis to
understand how AI-in-education policies are framed, how public concerns are
expressed, and where policy discourse and public sentiment align or diverge.

The methodology is structured as a connected analytical pipeline. Each layer
produces an output that informs the next layer, so the methods are not treated
as separate experiments. The pipeline begins with the construction of policy and
public sentiment corpora, prepares these texts for comparison, identifies
recurring themes, applies model-assisted classification, explores possible
policy--sentiment relationships, and finally evaluates the credibility of the
results.

The pipeline is organised into six layers:

\begin{enumerate}[label=\textbf{Layer \arabic*:}, leftmargin=*]

    \item \textbf{Corpus construction:} collects the policy documents and public
    sentiment sources that form the empirical basis of the study.

    \item \textbf{Preprocessing and segmentation:} transforms long and diverse
    documents into comparable text segments enriched with \textbf{metadata} such as
    source type, year, country or region, and thematic category.

    \item \textbf{Theme discovery:} uses topic modelling, clustering, and manual
    review to identify recurring themes across policy and sentiment texts.

    \item \textbf{Model-assisted classification:} uses SLMs and, where useful,
    LLM-supported coding to classify sentiment, stance, themes, barriers,
    enablers, and possible causal claims.

    \item \textbf{Causal NLP-inspired interpretation:} examines whether selected
    policy themes, such as privacy, regulation, teacher preparedness, equity, or
    AI literacy, are associated with public trust, concern, uncertainty, or
    perceived readiness.

    \item \textbf{Evaluation and triangulation:} checks whether the results are
    reliable, interpretable, and coherent across models, manual judgement,
    synthetic-data validation, source documents, and the literature review.

\end{enumerate}

This structure keeps the technical methods linked to the research question. The
corpus layer defines the evidence; preprocessing makes the evidence comparable;
theme discovery reveals candidate patterns; classification makes these patterns
more systematic; causal NLP-inspired interpretation explores possible
explanatory relationships; and evaluation checks whether the findings are
credible. The final aim is therefore not only to produce topics or model scores,
but to explain how the policy--practice gap in AI education governance is
expressed across policy discourse and public sentiment.

\subsection{Corpus construction and preprocessing}

The first two layers of the pipeline establish the empirical basis of the study.
The policy corpus will include national, European, international, and
institutional documents on AI in education. These documents will be analysed
with attention to themes such as:

\begin{itemize}[leftmargin=*]
    \item responsible and ethical AI;
    \item teacher preparedness and professional development;
    \item student AI literacy;
    \item equity, inclusion, and access;
    \item privacy, safety, and assessment;
    \item institutional capacity and implementation readiness.
\end{itemize}

The public sentiment corpus will include public-facing evidence on attitudes
towards AI in education and society, such as survey findings, reports,
consultation materials, and discourse sources where available. Each document or
text segment will be stored with metadata, including source type, country or
region where relevant, year, institutional origin, and thematic category.

\textbf{Long documents will be divided} into meaningful text segments so that different
sources can be compared without losing interpretive context. Text will then be
cleaned, deduplicated, segmented, and normalised where appropriate. Repeated
headers, references, and \textbf{non-substantive boilerplate} text will be removed. This
preprocessing stage will preserve the meaning of the original texts while making
them suitable for computational analysis.

\subsection{Exploratory topic and sentiment analysis}

After preprocessing, the next layer will identify recurrent themes and sentiment
patterns across the policy and public sentiment corpora. Topic modelling
methods, such as BERTopic or embedding-based clustering, will be used to
discover themes related to the topics mentioned above.

Topic modelling will be treated as a theme-discovery tool rather than a final
interpretation. \textbf{Generated topics will therefore be manually reviewed},
renamed, merged, or discarded where necessary. This stage will help identify whether
policy discourse emphasises the same issues that appear in public sentiment, or
whether some implementation barriers are underrepresented in formal policy language.

Sentiment and stance analysis will focus on categories that are meaningful for
AI education governance, including:

\begin{itemize}[leftmargin=*]
    \item trust and confidence;
    \item concern, fear, or resistance;
    \item optimism and perceived benefit;
    \item uncertainty or lack of readiness;
    \item perceived risk and perceived fairness.
\end{itemize}

Where off-the-shelf sentiment models are too general, small language models
(SLMs) can be prompted for narrower classification tasks. Candidate
SLMs may include transformer-based classifiers suitable for sentiment, stance,
theme, or barrier/enabler classification. Their role will be to support scalable
and repeatable labelling, not to replace human interpretation. Model outputs
\textbf{will be compared using evaluation measures} such as accuracy, precision, recall,
F1-score, and confusion matrices where labelled validation data is available.

\subsection{Use of SLMs, LLMs, and synthetic data}

The classification layer will use SLMs, LLM-assisted coding, and synthetic data
only where they strengthen the analysis. These tools will support theme
labelling, sentiment classification, stance detection, causal-claim detection,
and robustness testing, but they will not replace the source documents or human
interpretation.

Small language models will be used for focused and repeatable tasks. These may
include sentiment classification, stance detection, theme labelling, causal
claim detection, and barrier/enabler classification. SLMs are useful where the
task is narrow, well-defined, and requires consistent labelling across many text
segments.

Large language models (LLMs) will be used as supporting tools for tasks requiring
richer language understanding. These may include:

\begin{itemize}[leftmargin=*]
    \item summarising long policy sections;
    \item proposing initial qualitative codes;
    \item refining topic labels;
    \item identifying possible causal claims;
    \item generating controlled synthetic or counterfactual examples.
\end{itemize}

However, LLM outputs will not be treated as independent evidence. All
\textbf{LLM-assisted outputs will be checked against the source material}, and
important coding decisions will remain grounded in the corpus and the literature
review. \textbf{Prompt templates}, model versions, and \textbf{validation procedures}
will be documented to improve transparency.

Synthetic sentiment data will be used cautiously and only as a support
mechanism. It will not replace authentic public sentiment evidence. Its main
purposes will be:

\begin{itemize}[leftmargin=*]
    \item to address class imbalance;
    \item to create controlled examples for robustness testing;
    \item to support counterfactual testing of model behaviour;
    \item to test whether classifiers respond logically to changes in framing.
\end{itemize}

For example, synthetic variants may change a sentence from teacher support to
teacher replacement, or add conditions involving privacy, regulation, training,
or equity. These counterfactual examples will help test whether classification
models react consistently to changes in policy framing. They will not be
interpreted as real public opinion.

The faithfulness of synthetic data will be evaluated before it is used in any
analysis. Validation will include lexical comparison, semantic similarity using
sentence embeddings, topic-distribution comparison, sentiment-distribution
comparison, classifier-based real-versus-synthetic testing, and manual
inspection. The goal is not to prove that synthetic data is identical to real
data, but to determine whether it is safe and useful for augmentation,
robustness testing, and counterfactual model evaluation.

\subsection{Methodological roles, strengths, and limitations}

Table~\ref{tab:methodological_roles} summarises the role, advantages, and
limitations of the main methodological components. The purpose of the table is
to clarify which methods are central to the study and which ones act as
supporting tools.

\begin{table}[h]
\centering
\small
\caption{Methodological roles, strengths, and limitations}
\label{tab:methodological_roles}
\begin{tabular}{L{0.19\textwidth} L{0.27\textwidth} L{0.25\textwidth} L{0.22\textwidth}}
\hline
\textbf{Method} & \textbf{Role in the study} & \textbf{Main advantage} & \textbf{Main limitation} \\
\hline

Policy document analysis &
Identifies how AI in education is framed by institutions and governments. &
Connects the analysis to governance priorities such as ethics, equity, teacher
preparedness, and AI literacy. &
Policy texts may express aspirations rather than actual implementation capacity. \\[.15cm]

Public sentiment analysis &
Examines how people perceive AI, including trust, concern, optimism, uncertainty,
and resistance. &
Shows whether policy priorities correspond to public expectations and concerns. &
Sentiment can be context-dependent and may be simplified by generic categories. \\[.15cm]

Topic modelling and clustering &
Discovers recurrent themes across policy and sentiment texts. &
Helps compare large volumes of text and identify hidden patterns. &
Topics require manual interpretation and may not always map clearly to policy
concepts. \\[.15cm]

SLM-based classification &
Supports repeatable labelling of themes, sentiment, stance, barriers, and
enablers. &
Useful for focused, scalable, and consistent classification tasks. &
Requires validation and may struggle with complex or ambiguous policy language. \\[.15cm]

LLM-assisted analysis &
Supports summarisation, qualitative coding, theme refinement, and causal-claim
identification. &
Can help interpret complex texts and compare meanings across documents. &
Outputs must be checked and cannot be treated as independent evidence. \\[.15cm]

Synthetic data &
Supports class balancing, robustness testing, and controlled counterfactual
examples. &
Useful for testing whether models react logically to changes in framing. &
Must not replace real data and must be validated for faithfulness and relevance. \\[.15cm]

Causal NLP-inspired analysis &
Explores possible relationships between policy themes, implementation barriers,
and public sentiment. &
Moves the study from describing themes towards explaining possible links between
policy and practice. &
Findings should be presented as exploratory associations, not definitive causal
proof. \\[.15cm]

Validation metrics &
Assesses the reliability of models, labels, and synthetic data. &
Improves transparency and credibility of the computational pipeline. &
Metrics alone cannot replace qualitative judgement or theoretical interpretation. \\

\hline
\end{tabular}
\end{table}

\subsection{Causal NLP-inspired interpretation}

The causal NLP-inspired layer builds on the previous classification results. It
will not claim full causal inference. Instead, it will examine whether themes
identified in policy and sentiment texts appear to be associated with trust,
concern, uncertainty, perceived readiness, or implementation barriers.

For example, the study may examine whether:

\begin{itemize}[leftmargin=*]
    \item privacy-related discourse is associated with stronger concern;
    \item regulation language is linked to trust or reassurance;
    \item teacher-preparedness language relates to perceived readiness;
    \item innovation-focused language corresponds with optimism or anxiety;
    \item equity discourse is connected to concerns about unequal access.
\end{itemize}

Possible techniques include causal framing analysis, causal relation extraction,
and treatment-style comparisons between texts that mention or do not mention
selected themes. These findings will be presented as exploratory associations
and plausible explanatory patterns, not as definitive causal effects.

The interpretation stage will integrate the results across the different layers.
Policy themes will be compared with public sentiment themes to identify points
of alignment, divergence, and silence. Particular attention will be given to
teacher preparedness, student AI literacy, institutional readiness, equity,
privacy, public trust, and implementation barriers. The expected output is
therefore not only a set of model scores or topics, but an evidence-based
interpretation of where the AI education policy--practice gap appears, how it is
expressed, and what factors may help explain it.

\subsection{Evaluation of results}

The evaluation layer checks the reliability and meaning of the results produced
by the previous layers. It combines quantitative model evaluation, qualitative
manual review, synthetic-data validation, and triangulation with the source
documents and literature review.

For classification tasks, model outputs will be assessed using accuracy,
precision, recall, F1-score, and confusion matrices where labelled validation
data is available. For topic modelling and clustering, evaluation will focus on
topic coherence, interpretability, and manual review of whether the generated
topics correspond to meaningful policy or sentiment themes.

Synthetic data will be evaluated separately before being used for augmentation
or counterfactual testing. This will include lexical comparison, semantic
similarity, topic-distribution comparison, sentiment-distribution comparison,
classifier-based real-versus-synthetic testing, and manual inspection. Synthetic
data will only be used where it supports robustness testing or model evaluation;
it will not be treated as evidence of real public opinion.

The final findings will also be evaluated through triangulation. Results from
topic modelling, sentiment analysis, SLM/LLM-assisted classification, and causal
NLP-inspired interpretation will be compared against each other and against the
source documents. A finding will be treated as stronger when it is supported by
multiple layers of analysis and remains consistent with the literature review.
This evaluation layer is intended to ensure that the results are not only
technically measurable, but also meaningful for explaining the policy--practice
gap in AI education governance.
\end{document}